\section{Discussion}
\label{sec:discussion}

\subsection{Collapse and persistence from the same rules}

Simple local rules, none of which knows the size of the population, qualitatively reproduce much of what Calhoun
observed in Universe~25 when adults made non-breeding by crowding do not recover; when adults recover on their own
above a threshold rate, the colony resumes breeding and persists, oscillating rather than levelling off. The
pattern-oriented comparison of structures assigns the ingredients. The collapse needs crowding damage and long-lived
adults that do not recover; it survives every
change of structure that we applied, including the removal of the four social rules and of space, except adult recovery and
a short life span. The four social rules and space add specific patterns: the failure of the
late-born cohorts (social learning in a window), a persistent withdrawn class (retreat sites), the timing of the end
of natality (crowd aversion), and empty nest sites at a low peak with a prolonged stagnation (aggregation).

The mechanism connects old ideas of small-mammal ecology. Social stress from crowding depresses reproduction
\citep{Christian1950}, and an individual state carried across a delay makes density dependence
delayed \citep{GinzburgTaneyhill1994}. The model adds a developmental window in which
competence is learned socially \citep{FawcettFrankenhuis2015,Carcea2021} and shows that, with a long life span and a
non-breeding state that adults cannot leave, these close a loop: the adults whose competence crowding has removed
keep the crowding that prevents the next cohort from acquiring it. The loop acts on a specific class, the long-lived
non-breeders, which explains why indiscriminate mortality does not rescue the colony while their sustained removal
does (SM, Section~\ref{SM-supp:recovery}). The social component of the model is co-presence and
the transmission of a state, not social roles: dominance, male territoriality and reproductive skew, central to
Calhoun's account and to the social organisation of confined wild mice \citep{CrowcroftRowe1963}, are not modelled.

\subsection{Conditions of the mechanism}
\label{sec:conditions}

\paragraph{Irreversibility}
Density-dependent suppression of breeding in rodents is usually considered reversible. In house mice, urine of grouped females
delays puberty in young females \citep{Drickamer1977}, an effect that fades within ten days of ungrouping
\citep{CoppolaVandenbergh1985}. These signals act on the puberty of juvenile
females, the analogue of the juvenile window of the model, which is already reversible; direct evidence on the
recovery of adults that have stopped breeding is limited to deermice \citep{Terman1973}. With adult recovery
switched on, the same rules keep the colony alive, oscillating, instead of collapsing it
(Section~\ref{sec:res-recovery}). Near-irreversibility is
supported only by the transplants from Calhoun's colonies \citep{Marsden1972}, whose animals were several hundred
days old, when reproductive senescence of BALB/c females is already relevant; in \Cone\ the failure of the
transplants is reproduced by construction, not as independent evidence. The conclusion is therefore conditional:
simple rules give the collapse of Universe~25 if adults made non-breeding by crowding do not recover on their own
and live long (well mixed with purely social recovery no colony is rescued even at the juvenile rate, because no
competent demonstrators are left), and a persisting colony if they do recover. The threshold of adult recovery
is the quantity that an experiment on reversibility should estimate.

\paragraph{Scale}
Universe~25 was a single colony of one size. At fixed seeding density the model reproduces the patterns more often
in larger arenas, but the adult recovery that the collapse tolerates in space decreases with the arena, with no sign
of saturation between $L=15$ and 60 (Fig.~\ref{fig:recovery}c). Three sizes do not determine the limit of this
threshold for an infinite arena, which may be zero, so for any adult recovery above zero the spatial collapse may be
a finite-size effect; only the irreversible case is supported at every size simulated. Extinction itself is certain
for any closed colony with crowding damage, because crowding caps the number of animals that can breed and there is no
immigration; what distinguishes collapse from persistence is that the extinction time
of \Cone\ grows only logarithmically with the size of the colony (SM, Section~\ref{SM-supp:recovery}). The mechanism
is thus established at the scale of the experiment; a much larger closed colony would need an even less reversible
non-breeding state.

\paragraph{Timing}
The conjunction of all patterns is sensitive to when natality ends: with a background mortality of 0.003 per week, or
with initial ages spread over 4--80 weeks, it holds in 0.55 and 0.63 of the runs, only through the timing clause of
the end-of-natality criterion. The collapse itself is not sensitive: background mortality up to 0.05 per week does
not rescue \Cone\ (SM, Section~\ref{SM-supp:recovery}), a gestation-like refractory period of three weeks leaves it unchanged
(60 of 60 runs extinct; Section~\ref{sec:model}), and with spread initial ages the reversible null still falls to
0.23 of its peak and cycles (Section~\ref{sec:res-recovery}).

\subsection{Limits of the spatial layer}
\label{sec:spatial-limits}

The spatial patterns are obtained in a regime of mobility opposite to that of Universe~25. The ratio of the mixing
time to the deterioration time, a Damk\"ohler number, is about $1.7\times10^2$ in \Cone\ and at most $4\times10^{-3}$
in Universe~25, where mice crossed the enclosure daily (SM, Section~\ref{SM-supp:odd}): crossing the lattice by
diffusion takes about $1.7\times10^3$ weeks. In Universe~25 the classes arose among mobile animals through behaviour
(male territories, possession of nest boxes, aggression); in the model they arise from aggregation under low
mobility. The size of the withdrawn class is moreover set by a parameter: from week 30 the 270 retreat sites, with
540 places, are about 94\% occupied, so the class (416 adults at the peak in the colony of SM Fig.~\ref{SM-fig:snapshots}) is bounded
by $f_Rc_R|G|$. The proportion of withdrawn animals is thus an input of the model, not an output. The collapse, which
also occurs well mixed, does not depend on this regime; the classes do.

\subsection{Testable predictions}
\label{sec:predictions}

\emph{Adult recovery.} At the scale of the experiment ($L=30$), a closed colony collapses if fewer than about 1\%
of the adults of reproductive age breed again from 30 weeks after natality ends, and persists if more than about 9\%
do; the tolerated recovery decreases with colony size. It is measurable as the return of oestrous cycling, mating and maternal care in adults removed
from high density, as a function of age and exposure.
\emph{Sensitive window.} Cohorts reared at high density during the window fail even with competent demonstrators
(798 of 800 colonies); cross-fostering or translocation at different ages and densities would locate the window.
\emph{Removal of non-breeders.} Periodically removing a quarter to a half of the non-breeding adults rescues the
colony, while a single removal or indiscriminate mortality does not, a prediction for the management of captive
breeding colonies of social rodents (SM, Section~\ref{SM-supp:recovery}).
\emph{Retreat sites.} The withdrawn class scales with the number of small nest sites, $f_Rc_R|G|$; manipulating the
number of small nest boxes tests it.
These predictions concern \emph{Mus} in confinement. Elsewhere they depend on the reversibility of adult
suppression, which has been measured directly only in deermice \citep{Terman1973}.

\subsection{What is not reproduced}

The founding phase is reproduced only in an exploratory protocol with $a_{\max}=200$ weeks (about 3.8 years), a life
span implausible for \emph{Mus}. The simulated decline is a boom of one or two generations, not Calhoun's
multigenerational course: 69\% of the births are to the founders, and the net reproductive number falls from 3.8 to
0.40 in one generation. The persistent withdrawn class is made of refugees of the early boom, not of a late-born
cohort like the ``beautiful ones''; a variant that ties withdrawal to maternal socialisation produces a later-born
class at a demographic cost (O17p 0.705). Total infant mortality is absent, the female bias of the remnant holds in
fewer than half of the runs (0.45), and aggression is not modelled. The persisting colonies cycle with a period of
about 150 weeks instead of levelling off, and the litter-loss pattern of earlier colonies is reproduced only in part.
The same Calhoun patterns were used to choose the rules and to report them, so their fulfilment shows that they are
possible under local rules, not that they were predicted; only the comparison with earlier
confined colonies was held out, and it was defined after the production runs. A
quantitative model fitted to the census \citep{GwizdallaDus2026} reproduces the founding phase and the timing, which
ours does not, through switches on the total population; the two approaches agree that replicate variability is
large and that Calhoun's trajectory is one admissible outcome among many. We make no claim about human populations.
